\section{Tabel: Pengelompokan, Scope, dan Aksesibilitas}

Tabel yang lebih besar dan kompleks memerlukan pengelompokan baris agar struktur data mudah dipahami, baik oleh pengguna visual maupun teknologi bantu (assistive technology). HTML menyediakan tiga elemen pengelompokan: \code{<thead>} untuk baris header, \code{<tbody>} untuk baris data utama, dan \code{<tfoot>} untuk baris ringkasan atau total. Pengelompokan ini memberikan makna semantik tambahan dan memungkinkan browser menerapkan perilaku khusus---misalnya, \code{<thead>} dapat diulang saat tabel dicetak di beberapa halaman \cite{mdnHtmlIntro}.

Atribut \code{scope} pada elemen \code{<th>} menjelaskan cakupan header terhadap sel-sel data yang terkait. Nilai \code{scope="col"} menunjukkan bahwa header berlaku untuk seluruh kolom di bawahnya, sedangkan \code{scope="row"} menunjukkan header berlaku untuk seluruh baris di sebelah kanannya. Untuk tabel yang sangat kompleks---misalnya dengan header multi-level---dapat digunakan atribut \code{id} pada \code{<th>} dan \code{headers} pada \code{<td>} untuk membuat hubungan eksplisit antara sel data dan headernya \cite{mdnAccessibility}.

\begin{table}[htbp]
\centering
\begin{tabular}{|P{3.5cm}|P{4cm}|P{4.5cm}|}
\hline
\textbf{Mekanisme} & \textbf{Cara Kerja} & \textbf{Kapan Digunakan} \\
\hline
\code{scope="col"} & Header berlaku untuk kolom & Tabel sederhana, header di baris pertama \\
\hline
\code{scope="row"} & Header berlaku untuk baris & Tabel dengan header di kolom pertama \\
\hline
\code{scope="colgroup"} & Header untuk grup kolom & Tabel dengan header multi-level \\
\hline
\code{id} + \code{headers} & Hubungan eksplisit sel--header & Tabel sangat kompleks, header tidak beraturan \\
\hline
\end{tabular}
\caption{Perbandingan mekanisme penghubung header dan sel data}
\end{table}

Elemen \code{<caption>} memberikan judul deskriptif pada tabel dan sangat penting untuk aksesibilitas. Pembaca layar akan membacakan caption terlebih dahulu sehingga pengguna dapat memutuskan apakah tabel tersebut relevan sebelum mendengarkan seluruh isinya. Caption harus singkat namun informatif---menjelaskan apa yang disajikan tabel tanpa mengulangi data di dalamnya. Posisi visual caption dapat diatur dengan CSS menggunakan properti \code{caption-side} \cite{webdevAccessibility}.

\begin{figure}[htbp]
\centering
\resizebox{0.9\textwidth}{!}{%
\begin{tikzpicture}[node distance=1.2cm, transform shape, every node/.style={font=\footnotesize}]
  \node (table) [class, minimum width=10cm] {\code{<table>}};
  \node (caption) [object, below=0.5cm of table, minimum width=4cm] {\code{<caption>} Judul Tabel};
  \node (thead) [class, below=0.5cm of caption, minimum width=10cm, fill=blue!20] {\code{<thead>}};
  \node (th1) [object, below=0.4cm of thead, xshift=-2.5cm, minimum width=2cm] {\code{<th scope="col">}};
  \node (th2) [object, below=0.4cm of thead, xshift=0cm, minimum width=2cm] {\code{<th scope="col">}};
  \node (th3) [object, below=0.4cm of thead, xshift=2.5cm, minimum width=2cm] {\code{<th scope="col">}};
  \node (tbody) [class, below=1.2cm of thead, minimum width=10cm, fill=green!15] {\code{<tbody>}};
  \node (td1) [object, below=0.4cm of tbody, xshift=-2.5cm, minimum width=2cm] {\code{<td>} Data};
  \node (td2) [object, below=0.4cm of tbody, xshift=0cm, minimum width=2cm] {\code{<td>} Data};
  \node (td3) [object, below=0.4cm of tbody, xshift=2.5cm, minimum width=2cm] {\code{<td>} Data};
  \node (tfoot) [class, below=1.2cm of tbody, minimum width=10cm, fill=yellow!20] {\code{<tfoot>}};

  \draw [arrow] (table) -- (caption);
  \draw [arrow] (table.south) +(-4,0) |- (thead);
  \draw [arrow] (thead) -- (th1);
  \draw [arrow] (thead) -- (th2);
  \draw [arrow] (thead) -- (th3);
  \draw [arrow] (table.south) +(-4,0) |- (tbody);
  \draw [arrow] (tbody) -- (td1);
  \draw [arrow] (tbody) -- (td2);
  \draw [arrow] (tbody) -- (td3);
  \draw [arrow] (table.south) +(-4,0) |- (tfoot);
\end{tikzpicture}%
}
\caption{Struktur tabel HTML dengan pengelompokan thead, tbody, dan tfoot}
\end{figure}

Hindari penggunaan tabel bersarang (tabel di dalam tabel) karena sangat menyulitkan navigasi bagi pengguna pembaca layar. Jika data terlalu kompleks untuk satu tabel, pertimbangkan untuk memecahnya menjadi beberapa tabel yang lebih sederhana. Selalu pertimbangkan juga tampilan tabel di perangkat mobile---tabel yang lebar mungkin memerlukan scroll horizontal atau penyajian alternatif menggunakan CSS media queries \cite{mdnAccessibility}.

\begin{webcode}[language=HTML, caption={Tabel lengkap dengan thead, tbody, tfoot, scope, dan caption}]
<table>
  <caption>Rekap Nilai Mata Kuliah Pemrograman Web</caption>
  <thead>
    <tr>
      <th scope="col">NIM</th>
      <th scope="col">Nama</th>
      <th scope="col">UTS</th>
      <th scope="col">UAS</th>
    </tr>
  </thead>
  <tbody>
    <tr>
      <td>001</td> <td>Budi</td>
      <td>85</td>   <td>90</td>
    </tr>
    <tr>
      <td>002</td> <td>Siti</td>
      <td>78</td>   <td>88</td>
    </tr>
  </tbody>
  <tfoot>
    <tr>
      <th scope="row" colspan="2">Rata-rata</th>
      <td>81.5</td> <td>89</td>
    </tr>
  </tfoot>
</table>
\end{webcode}

\begin{catatan}
Untuk memastikan tabel Anda aksesibel, lakukan pengecekan berikut: (1) setiap tabel memiliki \code{<caption>}; (2) semua header menggunakan \code{<th>} dengan \code{scope} yang tepat; (3) hindari colspan/rowspan yang berlebihan; (4) uji tabel dengan pembaca layar atau navigasi keyboard saja. Alat seperti WAVE (\url{https://wave.webaim.org/}) dan aXe DevTools dapat membantu mengidentifikasi masalah aksesibilitas tabel secara otomatis.
\end{catatan}

